\Needspace{12\baselineskip}
\begin{center}
\textit{Historische Bemerkungen.}
\end{center}

\Needspace{8\baselineskip}
\Article{222}
Da mehrere der bisher behandelten Gegenstände auch von anderen Mathematikern untersucht wurden, können wir deren Verdienste nicht mit Schweigen übergehen. Der angesehene La Grange hat allgemeine Untersuchungen über die \textit{Äquivalenz der Formen} angestellt, \textit{Nouv. M\'em. de l'Ac. de Berlin,} 1773 \textit{p.} 263 und 1775 \textit{p.} 323 ff. Insbesondere zeigte er, dass es für jede gegebene Determinante eine endliche Anzahl von Formen gibt, sodass jede Form mit dieser Determinante einer davon äquivalent ist und sich somit alle Formen mit gegebener Determinante in Klassen einteilen lassen. Später entdeckte der berühmte Le Gendre mehrere elegante Eigenschaften dieser Klassifikation, größtenteils durch Induktion aus beobachteten Fällen; wir werden sie unten angeben und mit Beweisen begründen. Die Unterscheidung zwischen eigentlicher und uneigentlicher Äquivalenz, deren Nutzen besonders bei den feineren Untersuchungen sichtbar wird, hatte bisher jedoch niemand getroffen.

Das berühmte, in Artikel 216 ff. erläuterte Problem hat der angesehene La Grange zuerst vollständig\SourceLeaf{227} gelöst, \textit{Hist. de l'Ac. de Berlin,} 1767 \textit{p.} 165 und 1768 \textit{p.} 181 ff. Eine Lösung (allerdings eine weniger vollständige) steht auch in den schon öfter angeführten \textit{Suppl. ad Euleri Algebram}. Schon vorher hatte der angesehene Euler denselben Gegenstand in Angriff genommen, \textit{Comm. Petr. T.} VI \textit{p.} 175; \textit{Comm. Nov. T.} IX \textit{p.} 3; \textit{Ibid. T.} XVIII \textit{p.} 185 ff. Er beschränkte seine Untersuchung aber stets darauf, aus einer als bekannt vorausgesetzten Lösung andere abzuleiten. Außerdem können seine Methoden nur in wenigen Fällen \textit{alle} Lösungen liefern (siehe La Grange, \textit{Hist. de l'Ac. de Berlin} 1767, \textit{p.} 237). Da die letzte dieser drei Abhandlungen jüngeren Datums ist als La Granges Lösung, die das Problem in seiner ganzen Allgemeinheit erfasst und in dieser Hinsicht nichts zu wünschen übrig lässt, scheint Euler diese Lösung damals noch nicht gekannt zu haben (Band XVIII der Kommentare gehört zum Jahr 1773 und wurde 1774 veröffentlicht). Unsere Lösung beruht dagegen (wie alles andere, was wir bisher in diesem Abschnitt angegeben haben) auf ganz anderen Grundsätzen.

Was andere, etwa Diophant, Fermat usw., zu diesem Gegenstand beigetragen haben, betrifft sehr spezielle Fälle. Da wir das besonders erwähnenswert Erscheinende schon oben genannt haben, verzichten wir darauf, alles einzeln aufzuzählen.

\begin{center}
$\begin{array}{ccc}\ast&&\ast\\&\ast&\end{array}$
\end{center}

Was wir bisher über Formen zweiten Grades dargestellt haben, ist nur als erster Anfang dieser Theorie anzusehen. Als wir diese Untersuchung eingehender verfolgten, eröffnete sich uns ein sehr weites Feld, aus dem wir im Folgenden das besonders beachtenswert Erscheinende auswählen werden. Denn dieser Gegenstand ist so fruchtbar, dass wir der Kürze wegen vieles andere übergehen müssen, das uns bereits zu finden gelungen ist; weit mehr liegt ohne Zweifel noch verborgen und wartet auf neue Anstrengungen. Gleich zu Beginn dieser Untersuchungen ist zu bemerken, dass Formen mit der Determinante $0$ ausgeschlossen sind, sofern wir nichts anderes angeben.

\Needspace{12\baselineskip}
\begin{center}
\textit{WEITERE UNTERSUCHUNGEN ÜBER FORMEN.}
\end{center}

\Needspace{12\baselineskip}
\begin{center}
\textit{Einteilung der Formen mit gegebener Determinante in Klassen.}
\end{center}

\Needspace{8\baselineskip}
\Article{223}
Schon oben (Artikel 175, 195, 211) haben wir gezeigt: Für jede gegebene ganze Zahl $D$ (positiv oder negativ) lässt sich eine endliche Anzahl von Formen $F,F',F''$ usw. mit der Determinante $D$ angeben, sodass jede Form mit der Determinante\SourceLeaf{228} $D$ einer und nur einer davon eigentlich äquivalent ist. Alle Formen mit der Determinante $D$ (deren Anzahl unendlich ist) lassen sich daher nach diesen Formen \textit{klassifizieren}: Die Gesamtheit aller Formen, die $F$ eigentlich äquivalent sind, bildet die erste Klasse; die Formen, die $F'$ eigentlich äquivalent sind, bilden die zweite usw.

Aus jeder Klasse von Formen mit gegebener Determinante $D$ können wir eine Form wählen und als \textit{repräsentierende Form} der ganzen Klasse betrachten. An sich ist völlig beliebig, welche Form wir aus jeder Klasse nehmen. Dennoch sollten wir immer diejenige vorziehen, die die anderen an \textit{Einfachheit} zu übertreffen scheint. Die Einfachheit einer Form $(a,b,c)$ ist offenbar nach der Größe der Zahlen $a,b,c$ zu beurteilen. Mit Recht nennen wir eine Form $(a',b',c')$ weniger einfach als $(a,b,c)$, wenn $a'>a$, $b'>b$, $c'>c$ gelten. Damit ist die Sache aber noch nicht vollständig bestimmt: Es bleibt \textit{zum Beispiel} uns überlassen, welche der Formen $(17,0,-45)$, $(5,0,-153)$ wir für einfacher halten wollen. Meist wird es jedoch zweckmäßig sein, diese Regel einzuhalten:

I. Ist die Determinante $D$ negativ, so nehme man die in den einzelnen Klassen enthaltenen reduzierten Formen als repräsentierende Formen. Finden sich in derselben Klasse zwei reduzierte Formen (die dann entgegengesetzt sind, Artikel 172), so nehme man diejenige, deren mittleres Glied positiv ist.

II. Ist die Determinante $D$ positiv und keine Quadratzahl, so entwickle man die Periode einer reduzierten Form der vorgelegten Klasse. Darin werden entweder zwei ambige Formen vorkommen oder keine (Artikel 187).

1) Im ersten Fall seien die ambigen Formen $(A,B,C)$, $(A',B',C')$; die kleinsten Reste der Zahlen $B,B'$ nach den Moduln $A,A'$ seien jeweils $M,M'$ (sie können positiv genommen werden, sofern sie nicht $=0$ sind); schließlich sei $\frac{D-MM}{A}=N$, $\frac{D-M'M'}{A'}=N'$. Dann nehme man von den Formen $(A,M,-N)$, $(A',M',-N')$ diejenige als repräsentierende Form, die am einfachsten erscheint. Dabei ist eine Form vorzuziehen, deren mittleres Glied $=0$ ist. Ist das mittlere Glied in beiden oder in keiner $0$, so ziehe man die Form mit dem kleineren ersten Glied vor. Sind die ersten Glieder der Größe nach gleich, aber im Vorzeichen verschieden, so ist das positive Vorzeichen dem negativen vorzuziehen.

2) Gibt es in der ganzen Periode keine ambige Form, so wähle man aus allen Formen der Periode diejenige mit dem kleinsten ersten Glied, ohne Rücksicht auf das Vorzeichen\SourceLeaf{229}. Kommen in derselben Periode zwei Formen vor, in denen dasselbe erste Glied einmal positiv und einmal negativ ist, so ist die erste der zweiten vorzuziehen. Diese Form sei $(A,B,C)$. Daraus leite man wie im vorigen Fall eine andere Form $(A,M,-N)$ ab (indem man für $M$ den absolut kleinsten Rest von $B$ modulo $A$ nimmt und $N=\frac{D-MM}{A}$ setzt). Diese nehme man schließlich als Repräsentanten.

Sollte dasselbe kleinste erste Glied $A$ mehreren Formen der Periode gemeinsam sein, so behandle man alle diese Formen auf die vorgeschriebene Weise und nehme von den entstehenden Formen diejenige als repräsentierende Form, deren mittleres Glied möglichst klein wird.

{\rightskip=0pt plus 2em So haben wir \textit{zum Beispiel} für $D=305$ unter anderem diese Periode: $(17,4,-17)$, $(-17,13,8)$, $(8,11,-23)$, $(-23,12,7)$, $(7,16,-7)$, $(-7,12,23)$, $(23,11,-8)$, $(-8,13,17)$. Daraus wählen wir zunächst die Form $(7,16,-7)$ und leiten daraus dann die repräsentierende Form $(7,2,-43)$ ab.\par}

III. Ist die Determinante eine positive Quadratzahl $=kk$, so ermittle man die in der vorgelegten Klasse enthaltene reduzierte Form $(A,k,0)$. Ist $A<k$ oder $=k$, so nehme man diese selbst als repräsentierende Form. Ist dagegen $A>k$, so nehme man an ihrer Stelle die Form $(A-2k,k,0)$, deren erstes Glied negativ, aber kleiner als $k$ ist.

\textit{Beispiel.} So werden alle Formen mit der Determinante $-235$ in sechzehn Klassen eingeteilt. Ihre Repräsentanten sind $(1,0,235)$, $(2,1,118)$, $(4,1,59)$, $(4,-1,59)$, $(5,0,47)$, $(10,5,26)$, $(13,5,20)$, $(13,-5,20)$ und acht weitere, die sich von den vorigen nur durch die Vorzeichen der äußeren Glieder unterscheiden: $(-1,0,-235)$, $(-2,1,-118)$ usw.

Alle Formen mit der Determinante $79$ zerfallen in sechs Klassen, deren Repräsentanten $(1,0,-79)$, $(3.\,1,-26)$, $(3,-1,-26)$, $(-1,0.\,79)$, $(-3,1,26)$, $-3.\,-1,26)$ sind.

\Needspace{8\baselineskip}
\Article{224}
Durch diese Klassifikation werden also die eigentlich äquivalenten Formen vollständig von den übrigen getrennt. Zwei Formen derselben Determinante $D$ sind eigentlich äquivalent, wenn sie zur selben Klasse gehören. Jede durch die eine darstellbare Zahl lässt sich auch durch die andere darstellen. Lässt sich eine Zahl $M$ durch die erste Form mit teilerfremden Werten der Unbestimmten\SourceLeaf{230} darstellen, so ist dieselbe Zahl durch die andere Form ebenso darstellbar, und zwar so, dass beide Darstellungen zum selben Wert des Ausdrucks $\sqrt{D}\ (\mathrm{mod.}\ M)$ gehören. Gehören zwei Formen dagegen zu verschiedenen Klassen, so sind sie nicht eigentlich äquivalent. Von der Darstellbarkeit einer gegebenen Zahl durch die eine kann man nicht auf ihre Darstellbarkeit durch die andere schließen. Lässt sich vielmehr die Zahl $M$ durch die eine mit teilerfremden Werten der Unbestimmten darstellen, so sind wir sofort sicher, dass es keine ähnliche Darstellung derselben Zahl durch die andere Form gibt, die zum selben Wert des Ausdrucks $\sqrt{D}(\mathrm{mod.}\ M)$ gehört (siehe Artikel 167, 168).

Andererseits können zwei Formen $F,F'$ aus verschiedenen Klassen $K,K'$ durchaus uneigentlich äquivalent sein. Dann ist \textit{jede} Form der einen Klasse \textit{jeder} Form der anderen uneigentlich äquivalent. Jede Form aus $K$ hat ihre entgegengesetzte Form in $K'$, und die Klassen $K,K'$ selbst heißen \textit{entgegengesetzt}. So ist im ersten Beispiel des vorigen Artikels die dritte Klasse der Formen mit der Determinante $-235$ der vierten entgegengesetzt und die siebte der achten; im zweiten Beispiel die zweite Klasse der dritten und die fünfte der sechsten. Für beliebige zwei Formen aus entgegengesetzten Klassen gilt also: Jede Zahl $M$, die durch die eine darstellbar ist, lässt sich auch durch die andere darstellen. Geschieht dies bei der einen mit teilerfremden Werten der Unbestimmten, so ist es bei der anderen ebenso möglich, allerdings so, dass diese beiden Darstellungen zu entgegengesetzten Werten des Ausdrucks $\sqrt{D}\ (\mathrm{mod.}\ M)$ gehören. --- Die oben angegebenen Regeln zur Wahl der repräsentierenden Formen sind so eingerichtet, dass entgegengesetzte Klassen stets entgegengesetzte repräsentierende Formen erhalten.

Schließlich gibt es auch Klassen, die \textit{sich selbst entgegengesetzt} sind. Gehören nämlich eine Form und ihre entgegengesetzte zur selben Klasse, so sieht man leicht, dass alle Formen dieser Klasse einander sowohl eigentlich als auch uneigentlich äquivalent sind und ihre entgegengesetzten Formen in derselben Klasse haben. Jede Klasse, die eine ambige Form enthält, hat diese Eigenschaft. Umgekehrt enthält jede sich selbst entgegengesetzte Klasse notwendig eine ambige Form (Artikel 163, 165); deshalb heißt sie eine \textit{ambige Klasse}. So gibt es unter den Klassen der Formen mit der Determinante $-235$ acht ambige, deren Repräsentanten $(1,0,235)$, $(2,1,118)$, $(5,0,47)$, $(10,5,26)$, $(-1,0,-235)$, $(-2,1,-118)$, $(-5,0,-47)$, $(-10,5,-26)$ sind; unter den Klassen der Formen mit der Determinante $79$ zwei, mit den Repräsentanten $(1,0,-79)$, $(-1,0,79)$. --- Sind die repräsentierenden Formen nach unseren Regeln bestimmt, so lassen sich die ambigen Klassen leicht daran erkennen. Bei einer positiven\SourceLeaf{231} Nichtquadratzahl als Determinante erhält eine ambige Klasse sicher eine ambige repräsentierende Form (Artikel 194). Bei einer negativen Determinante ist die repräsentierende Form einer ambigen Klasse entweder selbst ambig oder hat gleiche äußere Glieder (Artikel 172). Bei einer positiven Quadratzahl als Determinante lässt sich schließlich nach Artikel 210 leicht entscheiden, ob die repräsentierende Form sich selbst uneigentlich äquivalent und somit die von ihr repräsentierte Klasse ambig ist.

\Needspace{8\baselineskip}
\Article{225}
Schon oben (Artikel 175) haben wir gezeigt, dass in einer Form $(a,b,c)$ mit negativer Determinante die äußeren Glieder dasselbe Vorzeichen haben, sowohl untereinander als auch wie die äußeren Glieder jeder äquivalenten Form. Sind $a,c$ positiv, so nennen wir die Form $(a,b,c)$ \textit{positiv}. Ebenso heißt die ganze Klasse, die $(a,b,c)$ enthält und nur aus positiven Formen besteht, eine \textit{positive Klasse}. Dagegen ist $(a,b,c)$ eine \textit{negative Form}, die in einer \textit{negativen Klasse} liegt, wenn $a,c$ negativ sind. Durch eine positive Form lassen sich keine negativen, durch eine negative Form keine positiven Zahlen darstellen. Ist die Form $(a,b,c)$ Repräsentant einer positiven Klasse, so ist die Form $(-a,b,-c)$ Repräsentant einer negativen Klasse. Daher gibt es gleich viele positive und negative Klassen; sobald die ersten angegeben sind, sind auch die zweiten bekannt. Bei Untersuchungen über Formen mit negativer Determinante genügt es deshalb meist, die positiven Klassen zu betrachten, deren Eigenschaften sich leicht auf die negativen Klassen übertragen.

Diese Unterscheidung gilt allerdings nur für Formen mit negativer Determinante. Durch Formen mit positiver Determinante lassen sich ohne Unterschied positive und negative Zahlen darstellen; nicht selten gehören sogar zwei Formen wie $(a,b,c)$, $(-a,b,-c)$ in diesem Fall zur selben Klasse.
